% BEGIN CR28_RADIAL_CONSTITUTIVE_EVALUATION
\section{Continuidad constitutiva y evaluación de la sección gravitatoria}
\label{cr28:sec:radial-evaluation}

La construcción longitudinal determina \(L\) antes de utilizar la
identificación gravitatoria. La composición de las reglas R1--R8 conserva
la sección de velocidad de \eqref{cr28:eq:same-velocity} y la acción
de retorno previamente generada. Con
\[
 L=\frac{5759}{23040}\alpha^{16}L_*,
 \qquad \hbar_{\rm ret}=\eta_{\rm ret}S_*,
 \qquad c_{\rm int}=c_*\widehat c,
\]
el resultado de \eqref{g26:eq:g-rigido-es} se expresa como
\begin{equation}
 G_{\rm ret}
 =\frac{c_*^3L_*^2}{S_*}
   \left(\frac{5759}{23040}\right)^2
   \frac{\alpha^{32}\widehat c^{\,3}}{\eta_{\rm ret}}
 =\frac{L^2}
 {\hbar_{\rm ret}(\mu_{\rm vac}\varepsilon_{\rm vac})^{3/2}}.
 \label{cr28:eq:g-constitutive-evaluation}
\end{equation}
Los canales son los mismos de \eqref{eq:vacuum-angular-channels}:
\[
 q_\pm=\exp\!\left[-\frac{\pi}{180}(A_{\rm deg}\pm C^*_{\rm deg})\right],
 \qquad r_\pm=\frac{1-q_\pm^{90}}{1-q_\pm^{120}},
 \qquad \widehat c=(r_+r_-)^{-1}.
\]
Las bases constitutivas satisfacen
\(\varepsilon=\varepsilon_*r_+^2\),
\(\mu=\mu_*r_-^2\) y
\(c_*=(\mu_*\varepsilon_*)^{-1/2}\). Su producto prueba
\(c_{\rm int}=(\mu\varepsilon)^{-1/2}\). Si se declara directamente
la coordenada dimensional de \(c_{\rm int}\), su base se recupera como
\(c_*=c_{\rm int}/\widehat c\). El coeficiente constitutivo no vuelve
a multiplicarse: ambas expresiones representan la misma sección.

En la representación \(L_*=1\,\mathrm m\),
\(S_*=10^{-34}\,\mathrm{J\,s}\) y
\(c_{\rm int}=299792458\,\mathrm{m\,s^{-1}}\), la evaluación de los
lectores y de la composición radial da
\begin{equation}
 \boxed{
 G_{\rm ret}
 =6.674299659414022706367776189\ldots\times10^{-11}
 \,\mathrm{m^3\,kg^{-1}\,s^{-2}}.}
 \label{cr28:eq:g-evaluation}
\end{equation}
El valor de contraste
\(G_{\rm CODATA\,2022}=6.67430(15)\times10^{-11}\,
\mathrm{m^3\,kg^{-1}\,s^{-2}}\) conserva su incertidumbre
experimental. En las mismas unidades, la diferencia de la evaluación
respecto del valor central es
\(-3.40585977\ldots\times10^{-18}\), aproximadamente
\(-0.00227057\) incertidumbres estándar. Las cifras de la evaluación
corresponden a los lectores y bases declarados; la incertidumbre pertenece
al resultado medido. Ninguna cifra de la referencia selecciona
\(N\), \(r_N\), la regla longitudinal o la norma radial.

\subsection{Realización cronológica del mismo radio construido}
\label{cr28:sec:radial-clock}
Con la duración elemental
\begin{equation}
 t_0=\frac{\pi L}{54c_{\rm int}},
 \qquad \ell_0=c_{\rm int}t_0=\frac{\pi L}{54},
 \label{cr28:eq:clock-from-length}
\end{equation}
el calendario de \(108\) posiciones satisface
\[
 T_{108}=108t_0=\frac{2\pi L}{c_{\rm int}},\qquad
 \omega_{108}=\frac{c_{\rm int}}{L},\qquad
 R_{108}=\frac{c_{\rm int}}{\omega_{108}}=L.
\]
Por sustitución,
\begin{equation}
 \left(\frac{54}{\pi}\right)^2
 \frac{c_{\rm int}^5t_0^2}{\hbar_a}
 =\frac{c_{\rm int}^3L^2}{\hbar_a}.
 \label{cr28:eq:circular-radial-agreement}
\end{equation}
Esta identidad sitúa el selector circular conservado a continuación en
la misma normalización longitudinal. Su teorema de unicidad conserva la
premisa geométrica que enuncia; la construcción anterior especifica además
la longitud utilizada en esa realización.

El carácter positivo \(Y=I\) da simultáneamente
\(R_X=\Lambda_X=L\) y \(M_X=\hbar_a/(c_{\rm int}L)\).
Para un carácter general \(Y>0\), las longitudes son recíprocas:
\(R_X=LY\), \(\Lambda_X=LY^{-1}\).
Su producto permanece \(L^2I\), sin que sus valores individuales deban
coincidir. Así, la identificación de las cuatro longitudes del caso
circular corresponde a un sector preciso de la prueba radial y no
suprime la dependencia del carácter.

La proporcionalidad \(R_X=(G/c_{\rm int}^{2})M_X\) es común a
todos los caracteres del dominio. Su conservación bajo cambio de base,
refinamiento, minimización con memoria, complemento de Schur y
normalización dinámica se ha demostrado en
\cref{g26:sec:rigidez-radial-es}. El portador pentacomponente y su
reducción a dos helicidades, desarrollados después, describen la
representación tensorial; no reemplazan la demostración del coeficiente
radial común.
% END CR28_RADIAL_CONSTITUTIVE_EVALUATION
